
This work began with a simple observation: practitioners universally default to LoRA rank $r=16$ without principled justification. The experiments demonstrate that optimal LoRA rank scales sub-linearly with model size, following $r_{\text{opt}} \propto N^{\alpha}$ where $\alpha = 0.82$ on SQuAD-v2 (95\% CI: [0.73, 0.92]).

However, two important caveats apply:

\begin{enumerate}
\item The scaling exponent is task-dependent: $\alpha \approx 0.30$ for multi-hop QA (HotpotQA), substantially different from single-hop QA.
\item The hypothesized mechanism (attention entropy) was refuted: larger models show lower, not higher, entropy.
\end{enumerate}

Larger models exhibit a phase transition in rank sensitivity---the penalty for suboptimal rank increases with scale---making principled rank selection increasingly important as models grow.

Future directions include:
\begin{enumerate}
\item Developing task complexity metrics that predict the scaling exponent $\alpha$
\item Testing whether optimal rank correlates with inverse attention entropy
\item Replicating experiments on additional architectures (Llama, Mistral)
\end{enumerate}

The present results provide initial evidence for scale-dependent LoRA rank optimization, with the caveat that task-specific calibration appears necessary.
